% Lesson 1 (AN-18): typing, word wrap, the status line, Backspace, clearing
% the screen without saving.  Replayed by test-pubs-iword.
% GATES: test-pubs-iword
\Gates{test-pubs-iword}
\Lesson{1}{Notes}
\LessonFeatures{Word Wrap; Status Line; Delete; Exit}
Instead of reaching for a pad when you need to leave a short note for a
colleague, type it.  In this lesson you type a note to Peter Gibbons about
the new cover sheets, correct the last word, and clear the screen.

\BeginReplay{l01}
\Task{Typing the Note}
As you type, press Shift for capital letters.  If you make a mistake, press
Backspace until the mistake is erased, then type the text again.
\begin{WPSteps}
\item Type the note below without pressing Enter at the end of a line:
\Type{Peter, the new cover sheets for the TPS reports are in the supply room. Please make sure that every report going out this week has one, and send me a copy of the memo when it is ready. Thanks, Bill}
\end{WPSteps}
Compare your screen with the one below.
\CompareScreen[rows=3]{l01-1}
Three things have happened as you typed.  First, the cursor has moved ahead
of the text: it marks where the next character goes.  Second, each time a
line filled, InitechWord moved the word that did not fit to the next line
by itself.  This is called \textbf{word wrap}.  Third, the status line at the
bottom of the screen has kept count.  \Scr{Doc 1} is the document, \Scr{Pg
1} the page, and \Scr{Ln 1.33"} the line the cursor is on, measured in
inches from the top of the page.  \Scr{Pos 7.6"} is how far the cursor is
from the left edge of the paper.

\Task{Correcting the Note}
Bill signs with a period.
\begin{WPSteps}
\item Press \PressKey{Backspace} four times to erase the name.\Keys{Backspace Backspace Backspace}
\item Type \TypeRun{Bill.}
\end{WPSteps}
\CompareScreen[rows=4]{l01-2}
With the period, the name no longer fits on the third line, so InitechWord
has moved it to the fourth.  The status line shows the cursor there:
\Scr{Ln 1.5"}, a sixth of an inch below the third line, and \Scr{Pos
1.5"}, one and a half inches from the left edge of the paper.

\Task{Clearing the Screen}
You do not need to keep the note.  To clear the screen without saving it,
\begin{WPSteps}
\item Press \Press{Exit}{F7}.  The status line asks
\Scr{Save Document? Yes (No)}.
\item Type \TypeName{n}.  The status line asks \Scr{Exit WP? No (Yes)}.
\CompareScreen[rows=4]{l01-3}
\item Type \TypeName{n} again to stay in InitechWord.  The screen clears.
\end{WPSteps}
\EndReplay
\begin{Aside}Typing y instead at Exit WP? would have closed InitechWord.  You
would start it again from the APPS folder, as described in the
Reference.\end{Aside}
In the next lesson you type a memo and save it, so that you can go on
working on it later.
